\chapter{Tắt Bớt Màn Hình Để Sống Lại}
\markboth{Tắt Bớt Màn Hình Để Sống Lại}{Tắt Bớt Màn Hình Để Sống Lại}

\begin{center}
	\textit{\color{bookprimary}``Cắt bớt màn hình không làm đời nghèo đi. Nó trả lại cho đời phần chú ý đã bị ăn trộm.''}

	\textit{\color{bookprimary}``Reducing screen time does not make life poorer. It returns the attention life has been robbed of.''}
\end{center}

\ngancach

\begin{tcolorbox}[colback=bookprimary!6,colframe=bookprimary!35,arc=5pt,title={\bfseries Góc Suy Nghĩ},fonttitle=\bfseries\color{bookprimary}]
Con người không thể sống mãi như một ngón tay biết lướt.
\textit{(A human being cannot live forever as a finger that knows only how to scroll.)}

Muốn lấy lại đời sống, phải lấy lại vài chỗ rất cụ thể.
\textit{(To reclaim life, one must reclaim a few very concrete places.)}

Bàn ăn.
\textit{(The dining table.)}

Giờ học.
\textit{(Study time.)}

Đường đi.
\textit{(The road.)}

Cuộc nói chuyện.
\textit{(The conversation.)}

Giấc ngủ.
\textit{(Sleep.)}

Và những lúc đầu óc đáng lẽ được yên.
\textit{(And the moments when the mind deserves quiet.)}
\end{tcolorbox}

\section{Một Kỷ Luật Nhỏ Nhưng Cứu Được Rất Nhiều}
\begin{truyen}{Một Kỷ Luật Nhỏ Nhưng Cứu Được Rất Nhiều}{Small Rules, Big Rescue}
\chuhoa{C}{ó lớp hỏi thầy}: ``Vậy phải bỏ điện thoại hẳn hả thầy?''
\textit{(A class asks the teacher: ``So do we have to give up phones completely?'')}

Thầy lắc đầu.
\textit{(The teacher shakes his head.)}

\teacherpair{Không. Vấn đề không phải là sống như người không có công nghệ. Vấn đề là đừng sống như nô lệ của nó. Em không cần vứt điện thoại. Em cần đặt lại địa vị của nó. Nó phải ở dưới đời sống của em, chứ không ngồi lên đầu em.}{``No. The issue is not to live as if technology does not exist. The issue is to stop living like its servant. You do not need to throw your phone away. You need to put it back in its proper rank. It should stay beneath your life, not sit on top of it.''}

Thầy viết lên bảng vài dòng rất cụ thể.
\textit{(The teacher writes a few very concrete lines on the board.)}

Không điện thoại trên bàn ăn.
\textit{(No phones on the dining table.)}

Không điện thoại khi người khác đang nói chuyện nghiêm túc với mình.
\textit{(No phones when someone is speaking seriously to you.)}

Không điện thoại khi chạy xe.
\textit{(No phones while driving or riding.)}

Không điện thoại trong những khoảng học cần tập trung sâu.
\textit{(No phones during study periods that require deep focus.)}

Không mang màn hình vào mọi hoạt động nghỉ ngơi của cơ thể.
\textit{(Do not bring screens into every form of bodily rest.)}

Rồi thầy nói chậm: ``Nghe nhỏ thôi, nhưng ai giữ được mấy điều này thật trong vài tháng sẽ thấy mình lấy lại rất nhiều phần đời tưởng đã mất.''
\textit{(Then the teacher says slowly: ``These rules sound small, but anyone who truly keeps them for a few months will feel many lost parts of life returning.'')}
\end{truyen}

\begin{tongketchuong}
Điện thoại không biến mất.
\textit{(The phone will not disappear.)}

Thuật toán cũng không tự thương con người hơn.
\textit{(Algorithms will not suddenly love human beings more either.)}

Điều còn lại là kỷ luật, là giới hạn, và là quyết định mỗi ngày xem mình muốn sống như một con người đầy đủ hay như một mắt nhìn luôn bị thuê lại bởi màn hình.
\textit{(What remains is discipline, limits, and the daily decision of whether one wants to live as a full human being or as a pair of eyes permanently rented out by the screen.)}
\end{tongketchuong}

\begin{goigiaovien}
Nếu dùng cuốn này để giáo dục học sinh, phụ huynh và chính mình, đừng kết bằng mắng.
\textit{(If this book is used to educate students, parents, and yourself, do not end with scolding.)}

Hãy kết bằng cam kết rất cụ thể.
\textit{(End with a very concrete commitment.)}

Một bữa ăn không màn hình.
\textit{(One screen-free meal.)}

Một giờ học không cầm máy.
\textit{(One study hour without holding the phone.)}

Một đoạn đường không nhìn xuống.
\textit{(One road segment without looking down.)}

Một cuộc nói chuyện không bị cắt bởi thông báo.
\textit{(One conversation not interrupted by notifications.)}

Giáo dục thật bắt đầu từ những thứ nhỏ nhưng được giữ nghiêm.
\textit{(Real education begins with small things kept seriously.)}
\end{goigiaovien}

\begin{tuhocnhanh}
1. Em sẽ lấy lại nơi nào trước: bàn ăn, giờ học, giấc ngủ hay đường đi? \textit{(Which place will you reclaim first: the dining table, study time, sleep, or the road?)}

2. Một quy tắc nhỏ nào em có thể giữ ngay từ hôm nay? \textit{(What small rule can you start keeping today?)}

3. Nếu bớt màn hình đi, em muốn phần đời nào của mình sống lại nhất? \textit{(If you reduce screen time, which part of your life do you most want to come back alive?)}
\end{tuhocnhanh}

\begin{tongketchuong}
Tắt bớt màn hình không phải là một phong trào đạo đức cho đẹp.
\textit{(Reducing screen time is not some pretty moral campaign.)}

Nhiều khi đó chỉ là cách cuối cùng để một người đòi lại mắt, đầu, tim và đời mình trước khi mọi thứ bị thuê mất quá sạch.
\textit{(Often it is simply the last way for a person to reclaim their eyes, mind, heart, and life before everything is rented out too completely.)}
\end{tongketchuong}

\ngancach